\subsection{格罗滕迪克群 R(A)}

设 A 是一个环。设 \(\mathcal{F}(A)\) 是有限生成的 A-模的类所成之集（VIII, 第 51 页）。有限长度的 A-模的类构成一个子集 \(\mathcal{L}\mathcal{F}(A)\)（\(\mathcal{F}(A)\) 的）；我们已看到 \(\mathcal{L}\mathcal{F}(A)\) 是模的类的一个遗传集。我们将与 \(\mathcal{L}\mathcal{F}(A)\) 相伴的格罗滕迪克群记作 R(A)，将有限长度的 A-模 E 的类在 R(A) 中的象记作 {[}E{]}。

第 3 小节的结果蕴含如下结论：

\begin{enumerate}
\def\labelenumi{\alph{enumi})}
\item
  设 \(\mathcal{S}(\mathrm{A})\) 是单 A-模的类所成之集。族 \(([S])_{S\in \mathcal{S}(A)}\) 是 \(\mathbf{Z}\)-模 \(R(A)\) 的一个基。
\item
  设 \(\mathrm{E}\) 与 \(\mathrm{F}\) 是有限长度的半单 A-模。则 \(\mathrm{E}\) 与 \(\mathrm{F}\) 同构当且仅当我们有 \([\mathrm{E}] = [\mathrm{F}]\)（在 \(\mathrm{R(A)}\) 中）。
\item
  设 E 是一个有限长度的 A-模，\((\mathrm{E}_{i})_{0\leqslant i\leqslant n}\) 是 E 的一个若尔当–赫尔德列。令 \(F=\bigoplus_{i=1}^{n}(\mathrm{E}_{i-1}/\mathrm{E}_{i})\)。则 F 是一个有限长度的半单 A-模，且在 R(A) 中我们有 \([E]=[F]\)。
\item
  设 \(\ell : \mathrm{R}(\mathrm{A}) \to \mathbf{Z}\) 是由对每个单 A-模 S 有 \(\ell([\mathrm{S}]) = 1\) 刻画的同态。于是对每个有限长度的 A-模 E，我们有 \(\ell([\mathrm{E}]) = \sum_{\mathrm{S} \in \mathcal{S}(\mathrm{A})} \ell_{\mathrm{S}}(\mathrm{E}) = \mathrm{long}_{\mathrm{A}}(\mathrm{E})\)。
\end{enumerate}

若 D 是一个体，则同态 \(\ell: \mathrm{R}(\mathrm{D}) \to \mathbf{Z}\) 是一个同构。

注. —— 1) 设 \(\mathcal{S}\mathcal{S}(\mathrm{A})\) 是有限长度的半单 A-模的类所成的遗传集。由 VIII, 第 190 页的命题 7，格罗滕迪克群 \(\mathrm{K}(\mathcal{S}\mathcal{S}(\mathrm{A}))\) 是一个自由 \(\mathbf{Z}\)-模，其元素 {[}S{]} \(_{\mathcal{S}\mathcal{S}(\mathrm{A})}\)（\(\mathrm{S} \in \mathcal{S}(\mathrm{A})\)）构成一个基。因此典范同态 \(\gamma_{\mathcal{L}\mathcal{F}(\mathrm{A}),\mathcal{S}\mathcal{S}(\mathrm{A})}\)（VIII, 第 186 页）是一个同构。

\begin{enumerate}
\def\labelenumi{\arabic{enumi})}
\setcounter{enumi}{1}
\item
  设 \(\mathrm{K}^{\prime}(\mathcal{L}\mathcal{F}(\mathrm{A}))\) 是 \(\mathcal{L}\mathcal{F}(\mathrm{A})\) 的直和格罗滕迪克群（VIII, 第 188 页）。克鲁尔–雷马克–施密特定理（VIII, 第 37 页）蕴含 \(\mathrm{K}^{\prime}(\mathcal{L}\mathcal{F}(\mathrm{A}))\) 是一个自由 Z-模，以有限长度的不可分解 A-模的类所成之集为基，而 \(\mathrm{K}(\mathcal{L}\mathcal{F}(\mathrm{A}))\) 以单 A-模的类所成之集为基。
\item
  设 E 是一个有限长度的 A-模。由 c)，存在一个有限长度的半单 A-模 \(E'\)，使得 \([E] = [E']\)；且由 b)，这样的模在相差一个同构的意义下是确定的；我们有时称之为 E 的半单化。
\end{enumerate}

命题 8. —— 设 A 是一个不是体的主理想整环，L 是它的分式域。存在一个同构 \(\varphi : \mathrm{R}(\mathrm{A}) \to \mathrm{L}^{*}/\mathrm{A}^{*}\)

使得

\[
\varphi ([ \mathrm{A} / a \mathrm{A} ]) = a \mathrm{A} ^ {*}\tag{9}
\]

对 A 中每个 \(a \neq 0\) 成立。

设 P 是由 A 的不可约元素组成的一个代表系（VII, §1, No.~3, 第 3 页）。A 的极大理想就是 \(p \in P\) 的理想 pA；因此每个单 A-模同构于唯一的模 A/pA。此外（VII, §1, No.~3, 第 4 页，定理 2），交换群 \(L^{*}/A^{*}\) 是自由的，且以族 \((pA^{*})_{p \in P}\) 为基。于是存在一个同构 \(\varphi\)（从 R(A) 到 \(L^{*}/A^{*}\)），它由 \(\varphi([A/pA]) = pA^{*}\)（对每个 \(p \in P\)）刻画。

设 A 中 \(a \neq 0\)。存在一个整数 \(r \geqslant 0\)、P 的元素 \(p_1, \ldots, p_r\) 与 A* 的一个元素 \(u\)，使得我们有 \(a = up_1 \cdots p_r\)。模 A/aA 具有由下式定义的合成列：

\[
\mathrm{E} _ {0} = \mathrm{A} / a \mathrm{A}, \quad \mathrm{E} _ {i} = (p _ {1} \dots p _ {i} \mathrm{A}) / a \mathrm{A} \quad (1 \leqslant i \leqslant r),
\]

且模 \(\mathrm{E}_{i-1}/\mathrm{E}_{i}=(p_{1}\cdots p_{i-1}\mathrm{A})/(p_{1}\cdots p_{i}\mathrm{A})\) 同构于 \(A/p_{i}A\)。因此我们有（VIII, 第 185 页，命题 3）

\[
\varphi ([ \mathrm{A} / a \mathrm{A} ]) = \prod_ {i = 1} ^ {r} \varphi ([ \mathrm{A} / p _ {i} \mathrm{A} ]) = p _ {1} \dots p _ {r} \mathrm{A} ^ {*} = a \mathrm{A} ^ {*}.
\]

注. —— 4) 我们保持命题 8 的假设与记号。设 E 是一个有限长度的 A-模，\(a_{1}\mathrm{A},\ldots ,a_{n}\mathrm{A}\) 是它的不变因子（VII, §4, No.~5, 第 20 页）。由于 E 同构于 \(\bigoplus_{i = 1}^{n}\mathrm{A} / a_i\mathrm{A}\)，我们有

\[
\varphi ([ \mathrm{E} ]) = \prod_ {i = 1} ^ {n} \varphi ([ \mathrm{A} / a _ {i} \mathrm{A} ]) = a _ {1} \dots a _ {n} \mathrm{A} ^ {*}.\tag{10}
\]

\begin{enumerate}
\def\labelenumi{\arabic{enumi})}
\setcounter{enumi}{4}
\tightlist
\item
  *设 A 是一个不是体的戴德金环（Comm. Alg., VIII, §2, No.~1）。按与命题 8 中同样的推理，我们可以证明存在一个从 R(A) 到 A 的分式理想群的同构 \(\varphi\)，它由 \(\varphi([A/\mathfrak{a}]) = \mathfrak{a}\)（对 A 的每个非零理想 \(\mathfrak{a}\)）刻画。*
\end{enumerate}

命题 8 将在例如下面两种情形中被用到：

\begin{enumerate}
\def\labelenumi{\alph{enumi})}
\tightlist
\item
  假设我们有 A = Z。有限长度的 Z-模就是有限交换群。由于 \(Q^{*}\) 是 \(Z^{*} = \{1, -1\}\) 与 \(Q_{+}^{*}\) 的直积，我们可以由命题 8 导出一个同构 \(\varphi'\)（从 \(\mathrm{R}(\mathbf{Z})\) 到 \(Q_{+}^{*}\)），它由
\end{enumerate}

\[
\varphi^ {\prime} ([ \mathrm{G} ]) = \operatorname{Card} (\mathrm{G})
\]

（对每个有限交换群 G）给出。

\begin{enumerate}
\def\labelenumi{\alph{enumi})}
\setcounter{enumi}{1}
\tightlist
\item
  假设 A 是交换域 K 上一个未定元 T 的多项式环 K{[}T{]}。设 E 是 K 上一个有限维向量空间，u 是 E 的一个自同态。如同 VII, §5, No.~1, 第 29 页中那样，用 \(E_{u}\) 表示以 E 为底层加群、以 \((p, x) \mapsto p(u)x\) 为外部规律的 A-模。A-模 \(E_{u}\) 具有有限长度。反之，每个单 A-模在 K 上是有限维的（VII, §4, No.~8, 第 25 页，注 4）。因此，每个有限长度的 A-模在 K 上是有限维的，从而形如 \(E_{u}\)。此外（VII, §5, No.~3, 第 32 页，推论 1），\(E_{u}\) 的诸不变因子之积等于 u 的特征多项式 \(\chi_{u}\)。于是，命题 8 给出一个同构
\end{enumerate}

\[
\varphi : \mathrm{R} (\mathrm{K} [ \mathrm{T} ]) \to \mathrm{K} (\mathrm{T}) ^ {*} / \mathrm{K} ^ {*}
\]

它由 \(\varphi([\mathrm{E}_u]) = \chi_u\mathrm{K}^*\) 刻画（参见公式 (10)）。

